\documentclass[11pt]{article}
\usepackage[a4paper,margin=25mm]{geometry}
\usepackage{amsmath,amssymb,amsthm}
\usepackage{longtable,booktabs,array}
\usepackage[hidelinks]{hyperref}
\hypersetup{pdftitle={E22567506\_5: 2x\textasciicircum 4+xy+x+y\textasciicircum 3-y+1=0 has no integer solutions}}
\newtheorem{theorem}{Theorem}
\newtheorem{lemma}[theorem]{Lemma}
% keep the space around binary operators in inline formulas
\medmuskip=4mu plus 2mu\relax
\emergencystretch=1.5em
\tolerance=400
\allowdisplaybreaks
\newcommand{\Z}{\mathbb{Z}}
\newcommand{\Q}{\mathbb{Q}}
\newcommand{\F}{\mathbb{F}}
% Legendre and Jacobi symbols: one fixed size in running text, one in displays
\newcommand{\leg}[2]{\mathchoice{\biggl(\dfrac{#1}{#2}\biggr)}{(#1/#2)}{(#1/#2)}{(#1/#2)}}
\title{\textbf{E22567506\_5}\\[4pt] {\Large The equation $2x^{4}+xy+x+y^{3}-y+1=0$ has no integer solutions}}
\author{}
\date{}
\begin{document}
\maketitle
\vspace{-8ex}
\begin{center}
Size $h=49$, length $l=12$
\end{center}

\begin{theorem}\label{thm:main}
The equation
\begin{equation}\label{eq:main}
2x^{4}+xy+x+y^{3}-y+1=0
\end{equation}
has no integer solutions.
\end{theorem}

The proof uses the cubic Hilbert reciprocity law over the field $K=\Q(\omega)$, where $\omega^2+\omega+1=0$. It uses 2 graphs. For the $g$-th graph, we attach to a hypothetical integer solution the values of auxiliary polynomials $H^{(g)}_i$ with coefficients in $\Z[\omega]$, constants $c^{(g)}_i\in\Z[\omega]$, and a sum $B^{(g)}$ of cubic Hilbert symbols of these values, recorded by weighted edges. By the reciprocity law, the local contributions $B^{(g)}_p$ of all primes $p$ add up to $0$ modulo $3$, for every $g$. We show that $B^{(g)}_p=0$ for $p\notin S$, compute the vectors $(B^{(g)}_p)_g$ for the primes $p\in S$, and obtain a contradiction. Two standard facts from class field theory are quoted without proof: the explicit formula for the cubic Hilbert symbol at the places not above $3$, and the Hilbert reciprocity law. Apart from these, the proof uses polynomial identities, which can be checked by expanding both sides, and finite computations with residues.


\section{\texorpdfstring{Cubic Hilbert symbols over $\Q(\omega)$}{Cubic Hilbert symbols over Q(w)}}\label{sec:symbols}

\paragraph{The field.} Let $K=\Q(\omega)$, where $\omega^2+\omega+1=0$. Its ring of integers is $\Z[\omega]$, and every element of $K$ can be written uniquely as $A+B\omega$ with $A,B\in\Q$. We write
\[
\langle A,B\rangle=A+B\omega,
\]
also when $A$ and $B$ are polynomials in $x$ and $y$. Then $\langle A,B\rangle\langle C,D\rangle=\langle AC-BD,\ AD+BC-BD\rangle$, and the norm is $N(\langle A,B\rangle)=A^2-AB+B^2$. In particular, a value $\langle A,B\rangle$ with $A,B\in\Z$ is zero if and only if $A=B=0$.

\paragraph{Places.} The finite places of $K$ correspond to the non-zero prime ideals of $\Z[\omega]$. A prime $p\equiv2\pmod3$ is inert: there is one place above $p$, with uniformizer $p$ and residue field $\Z[\omega]/p\Z[\omega]$ with $p^2$ elements. If $p\equiv1\pmod3$, the polynomial $t^2+t+1$ has two roots $r$ modulo $p$, and for each of them there is a place $v_r$ above $p$, at which $\omega$ is the root of $t^2+t+1$ in $\Z_p$ congruent to $r$ modulo $p$; the completion is $\Q_p$, the uniformizer is $p$, the residue field is $\F_p$, and $\omega$ reduces to $r$. The prime $3$ is ramified: $\pi=1-\omega$ generates the unique prime ideal above $3$, $\pi^2=-3\omega$, and $\Z[\omega]/\pi\Z[\omega]=\F_3$, where $\langle A,B\rangle$ reduces to $A+B$. The only infinite place of $K$ is complex. For a finite place $v$ we write $K_v$ for the completion and $v(a)$ for the normalized valuation, and we call $a$ a \emph{unit} at $v$ if $v(a)=0$.

\paragraph{Hilbert symbols.} For a place $v$ of $K$ and $a,b\in K_v^\times$, the cubic Hilbert symbol $(a,b)_v$ is a cube root of unity, see \cite[Chapter~V, \S3]{N}. We write $(a,b)_v=\omega^{\beta_v(a,b)}$ with $\beta_v(a,b)\in\Z/3\Z$, and we use the following three facts \cite[Chapter~V, \S3 and Chapter~VI, \S8]{N}.
\begin{itemize}
\item[(H1)] $\beta_v(a,bc)=\beta_v(a,b)+\beta_v(a,c)$, $\beta_v(a,b)=-\beta_v(b,a)$, and $\beta_v(a,c^3)=0$. In particular, $\beta_v(a,b)$ depends only on the classes of $a$ and $b$ in $K_v^\times/K_v^{\times3}$.
\item[(H2)] Let $v$ be a finite place not above $3$, let $q$ be the number of elements of its residue field, and let $\varpi$ be a uniformizer at $v$. Write $a=\varpi^eu$ and $b=\varpi^fw$ with units $u$ and $w$, and define $\chi(u)\in\Z/3\Z$ by $\bar u^{(q-1)/3}=\bar\omega^{\chi(u)}$, where the bar denotes the reduction to the residue field. Then $\beta_v(a,b)=e\,\chi(w)-f\,\chi(u)$. In particular, $\beta_v(a,b)=0$ if $a$ and $b$ are both units.
\item[(H3)] (Hilbert reciprocity.) For $a,b\in K^\times$, $\beta_v(a,b)=0$ for all but finitely many places $v$, and $\sum_v\beta_v(a,b)=0$, where the sum is over all places of $K$. At the complex place, $\beta_v(a,b)=0$.
\end{itemize}
The general formula in (H2) contains an additional term $ef\chi(-1)$, which vanishes because $-1=(-1)^3$ is a cube. Formula (H2) fixes the normalization of the symbols; with the opposite sign convention all values $\beta_v$ change sign, and the proof below is the same.

\paragraph{Classes at the places not above $3$.} Let $p\neq3$ be a prime and $v$ a place above $p$. For $a=p^eu\in K_v^\times$ with a unit $u$, put $c_v(a)=(e\bmod3,\ \chi(u))\in\F_3^2$. By (H2) and (H1), $c_v$ is additive, $c_v(a)$ determines the class of $a$ modulo cubes, and $\beta_v(a,b)=e\chi(w)-f\chi(u)$ for $c_v(a)=(e,\chi(u))$ and $c_v(b)=(f,\chi(w))$. We put $c_p(a)=c_v(a)\in\F_3^2$ if $p\equiv2\pmod3$, and $c_p(a)=(c_{v_r}(a),c_{v_{r'}}(a))\in\F_3^4$ if $p\equiv1\pmod3$, where $r<r'$ are the roots of $t^2+t+1$ in $\{1,\dots,p-1\}$; accordingly, $\beta_p$ is the sum of the forms $\beta_v$ over the places $v$ above $p$. For $p\equiv2\pmod3$ and a rational integer $u$ prime to $p$, we have $\bar u\in\F_p$, so that $\bar u^{(p^2-1)/3}=(\bar u^{p-1})^{(p+1)/3}=1$ and $\chi(u)=0$.

\paragraph{The place above $3$.}
\begin{lemma}\label{lem:M3}
Let $v$ be the place above $3$. Every non-zero element of $K_v$ is equal to $\pi^a\omega^b2^c(1+3\omega)^d$ times a cube, for unique $a,b,c,d\in\Z/3\Z$, and every unit congruent to $1$ modulo $9$ is a cube. Put $c_3(A)=(a,b,c,d)\in\F_3^4$. Then $c_3$ is additive, and $\beta_v(A,B)=c_3(A)^{\mathsf T}M_3\,c_3(B)$, where
\[
M_3=\begin{pmatrix}0&0&2&0\\0&0&1&2\\1&2&0&0\\0&1&0&0\end{pmatrix}\pmod3.
\]
\end{lemma}

\begin{proof}
For $\alpha\in\Z[\omega]$, the polynomial $g(s)=s+3s^2+3s^3-\alpha$ satisfies $g'(s)\equiv1\pmod\pi$, so by Hensel's lemma it has a root $s$ in the completion of $\Z[\omega]$, and then $(1+3s)^3=1+9(s+3s^2+3s^3)=1+9\alpha$. Hence every unit congruent to $1$ modulo $9$ is a cube in $K_v$, and the class of a unit modulo cubes depends only on its residue modulo $9$. There are $54$ units in $\Z[\omega]/(9)$, and a direct enumeration shows that the $27$ products $\omega^b2^c(1+3\omega)^d$ with $b,c,d\in\{0,1,2\}$ are pairwise distinct modulo $9$, also after multiplication by $-1=(-1)^3$. Hence they represent all $27$ classes of units modulo cubes. Since $\pi^3$ is a cube times a unit, the exponent of $\pi$ is the valuation modulo $3$, and the first statement follows.

By (H1), it remains to compute $\beta_v$ on the $16$ ordered pairs of the elements $\pi$, $\omega$, $2$ and $1+3\omega$. Their norms are $3$, $1$, $4$ and $7$, hence, for each such pair, all the symbols at the places not above $2$, $3$ and $7$ vanish by (H2), and, by (H3), $\beta_v$ at the place above $3$ is minus the sum of the values at the place above $2$ and at the two places above $7$. These places correspond to $\bar\omega=2$ and $\bar\omega=4$ in $\F_7$. The classes $(e,\chi)$ of the four elements at these places are
\[
\begin{array}{c|ccc}
 & 2 & \bar\omega=2 & \bar\omega=4\\\hline
\pi & (0,2) & (0,0) & (0,2)\\
\omega & (0,1) & (0,2) & (0,2)\\
2 & (1,0) & (0,2) & (0,1)\\
1+3\omega & (0,2) & (1,0) & (0,0)
\end{array}
\]
where, at the places above $7$, $\chi(u)$ is defined by $\bar u^2=\bar\omega^{\chi(u)}$ in $\F_7$, and $7$ is a uniformizer. Formula (H2) and the sum over these three places give $M_3$.
\end{proof}

We write $\beta_3$ for the form of Lemma~\ref{lem:M3}. Since $3=-\omega^2\pi^2$ and $4=2^2$, we have $c_3(3)=(2,2,0,0)$ and $c_3(4)=(0,0,2,0)$.

\paragraph{Residue classes.} The local computations use the following description of the classes of the values of a polynomial on a residue class.

\begin{lemma}\label{lem:residue}
Let $p$ be a prime, $k\geq1$, and let $z\in\Z[\omega]$ be non-zero with $z\equiv\zeta\pmod{p^k}$, where $\zeta=\langle A,B\rangle$ with $0\leq A,B<p^k$. Then $c_p(z)\in b+D$ for the vector $b$ and the subspace $D$ given as follows.
\begin{itemize}
\item[(i)] $p\equiv2\pmod3$. If $\zeta\neq0$, then $b=c_p(\zeta)$ and $D=0$. If $\zeta=0$, then $b=0$, and $D$ is $\F_3^2$, or the span of $(1,0)$ if $z\in\Z$.
\item[(ii)] $p\equiv1\pmod3$. For each root $r$, let $\rho$ be the root of $t^2+t+1$ modulo $p^k$ congruent to $r$. If $A+B\rho\not\equiv0\pmod{p^k}$, the two coordinates of $c_{v_r}(z)$ are those of $c_{v_r}(\zeta)$; otherwise these two coordinates are arbitrary.
\item[(iii)] $p=3$, $z\notin\Z$. If $\zeta\neq0$, let $a=v(\zeta)<2k$ be its valuation at $\pi$ and $m=2k-a$. Then $b=c_3(\zeta)$, and $D=0$ if $m\geq4$, while $D=D_m$ if $m\leq3$, where $D_1$ is spanned by $(0,1,0,0)$, $(0,0,1,0)$ and $(0,0,0,1)$, $D_2$ by $(0,0,1,0)$ and $(0,0,0,1)$, and $D_3$ by $(0,0,1,1)$. If $\zeta=0$, then $b=0$ and $D=\F_3^4$.
\item[(iv)] $p=3$, $z\in\Z$, so that $B=0$. If $3^k\nmid A$, let $j=v_3(A)$. Then $b=c_3(A)$ and $D=0$ if $k-j\geq2$, while $b=j\,c_3(3)$ and $D$ is spanned by $c_3(4)$ if $k-j=1$. If $3^k\mid A$, then $b=0$ and $D$ is spanned by $c_3(3)$ and $c_3(4)$.
\end{itemize}
\end{lemma}

\begin{proof}
(i) If $\zeta\neq0$, then $e=v(z)=v(\zeta)<k$ and $z/p^e\equiv\zeta/p^e\pmod p$, so that $\chi$ is determined. If $z\in\Z$, its unit part is a rational integer, and $\chi=0$. (ii) At $v_r$, $z$ and $\zeta$ are the $p$-adic numbers $A'+B'\omega$ and $A+B\omega$ with $\omega\equiv\rho\pmod{p^k}$, and the argument of (i) applies. (iii) If $\zeta\neq0$, then $v(z)=a$, and $z/\pi^a\equiv\zeta/\pi^a\pmod{\pi^m}$. By Lemma~\ref{lem:M3}, the class of a unit depends only on its residue modulo $9=\omega\pi^4$; the units congruent to $1$ modulo $\pi^m$ form a subgroup, and a direct enumeration of their residues modulo $9$ shows that their classes form $D_m$. (iv) If $3^k\nmid A$, then $z=3^ju$ with a rational unit $u\equiv A/3^j\pmod{3^{k-j}}$. The class of a rational unit depends only on its residue modulo $9$, the units $\pm1$ are cubes, and the rational units congruent to $1$ modulo $3$ are $4^s$ times a unit congruent to $1$ modulo $9$; hence the classes are as stated, and if $3^k\mid A$ the valuation $j$ is also unknown.
\end{proof}

\begin{lemma}\label{lem:envelope}
Let $p$ be a prime, and consider vertices $i$ with vectors $b_i$, subspaces $D_i$ and vectors $\gamma_i$, and a set of edges $(i,j)$ with $i<j$ and weights $w_{ij}\in\{1,2\}$. Put $\rho_{ij}=w_{ij}$ and $\rho_{ji}=-w_{ij}$ for the edges, and $\sigma_i=\gamma_i+\sum_j\rho_{ij}b_j$. Suppose that $\beta_p(d,\sigma_i)=0$ for every $i$ and every $d\in D_i$, and that $\beta_p(d,d')=0$ for every edge $(i,j)$ and all $d\in D_i$, $d'\in D_j$. Then
\[
\sum_{(i,j)}w_{ij}\beta_p(c_i,c_j)+\sum_i\beta_p(c_i,\gamma_i)=\sum_{(i,j)}w_{ij}\beta_p(b_i,b_j)+\sum_i\beta_p(b_i,\gamma_i)
\]
whenever $c_i\in b_i+D_i$ for every $i$. By bilinearity, it suffices to check the hypotheses for $d$ and $d'$ in spanning sets of $D_i$ and $D_j$.
\end{lemma}

\begin{proof}
Write $c_i=b_i+d_i$ with $d_i\in D_i$. As $\beta_p$ is bilinear and $\beta_p(u,v)=-\beta_p(v,u)$, the left side equals the right side plus $\sum_i\beta_p(d_i,\sigma_i)+\sum_{(i,j)}w_{ij}\beta_p(d_i,d_j)$, and both sums vanish.
\end{proof}
\section{The graphs}

Put $F(x,y)=2x^{4}+xy+x+y^{3}-y+1$, so that \eqref{eq:main} is the equation $F(x,y)=0$. There are 2 certificates. The $g$-th certificate consists of polynomials $H^{(g)}_i$ with coefficients in $\Z[\omega]$ and constants $c^{(g)}_i\in\Z[\omega]$, listed in Table~\ref{tab:vert1} and Table~\ref{tab:vert2}, and of \emph{edges} $(i,j)$ with $i<j$ and weights $w^{(g)}_{ij}\in\{1,2\}$. For an edge we put $\rho^{(g)}_{ij}=w^{(g)}_{ij}$ and $\rho^{(g)}_{ji}=-w^{(g)}_{ij}\bmod3$, and $\rho^{(g)}_{ij}=0$ otherwise. In the rest of the proof we usually suppress the index $g$.
The edges $(i,j;w_{ij})$ of the first certificate are $(0,1;2)$, $(0,4;1)$, $(0,5;1)$, $(0,7;1)$, $(0,10;1)$, $(0,13;2)$, $(0,15;1)$, $(0,16;2)$, $(0,18;1)$, $(0,27;2)$, $(1,2;2)$, $(1,3;2)$, $(1,5;1)$, $(1,8;1)$, $(1,11;1)$, $(1,12;2)$, $(1,15;1)$, $(1,23;1)$, $(2,5;2)$, $(2,9;2)$, $(2,16;2)$, $(2,17;2)$, $(3,12;2)$, $(3,16;1)$, $(3,18;2)$, $(3,19;1)$, $(3,23;1)$, $(4,18;2)$, $(4,23;1)$, $(4,24;1)$, $(4,27;2)$, $(5,9;1)$, $(5,19;1)$, $(5,23;2)$, $(6,9;2)$, $(6,10;1)$, $(6,11;2)$, $(6,18;1)$, $(6,21;1)$, $(6,28;2)$, $(7,9;2)$, $(7,14;1)$, $(7,19;1)$, $(8,9;1)$, $(8,15;2)$, $(8,16;1)$, $(9,14;2)$, $(9,20;2)$, $(9,22;1)$, $(10,20;1)$, $(10,22;1)$, $(10,25;2)$, $(11,15;2)$, $(11,16;1)$, $(11,17;2)$, $(11,22;1)$, $(11,27;2)$, $(12,13;1)$, $(12,14;1)$, $(12,28;2)$, $(13,14;1)$, $(13,15;2)$, $(13,16;2)$, $(13,21;2)$, $(13,27;2)$, $(14,15;1)$, $(15,19;1)$, $(17,21;2)$, $(17,27;1)$, $(18,24;2)$, $(18,26;1)$, $(19,24;1)$, $(19,26;2)$, $(20,21;1)$, $(21,25;1)$, $(22,28;2)$, $(23,26;1)$, $(24,26;1)$, $(24,27;1)$, $(26,28;2)$.
The edges $(i,j;w_{ij})$ of the second certificate are $(0,15;1)$, $(0,26;1)$, $(1,6;2)$, $(1,11;2)$, $(1,12;1)$, $(1,22;2)$, $(2,5;2)$, $(2,6;1)$, $(2,7;2)$, $(2,8;1)$, $(2,10;1)$, $(2,11;1)$, $(2,13;1)$, $(2,17;1)$, $(2,18;1)$, $(2,19;2)$, $(2,21;1)$, $(2,26;1)$, $(2,29;2)$, $(3,4;2)$, $(3,8;1)$, $(3,16;1)$, $(3,21;2)$, $(3,23;1)$, $(3,25;2)$, $(4,5;1)$, $(4,17;1)$, $(4,20;1)$, $(4,21;1)$, $(4,22;1)$, $(4,23;1)$, $(4,29;2)$, $(5,10;1)$, $(5,14;2)$, $(5,16;1)$, $(5,18;1)$, $(5,19;1)$, $(5,23;2)$, $(5,24;1)$, $(5,25;1)$, $(5,31;1)$, $(6,12;1)$, $(6,17;2)$, $(6,19;1)$, $(6,20;1)$, $(6,27;2)$, $(6,30;2)$, $(7,8;1)$, $(7,9;2)$, $(7,14;1)$, $(7,16;2)$, $(7,17;1)$, $(7,19;1)$, $(7,22;2)$, $(7,25;1)$, $(7,29;1)$, $(8,9;1)$, $(8,13;1)$, $(8,15;2)$, $(8,17;2)$, $(8,19;2)$, $(8,21;2)$, $(8,22;2)$, $(8,29;1)$, $(9,21;1)$, $(10,15;2)$, $(10,27;1)$, $(10,30;2)$, $(11,18;1)$, $(11,20;2)$, $(11,21;2)$, $(11,24;2)$, $(11,27;1)$, $(12,23;1)$, $(12,25;2)$, $(12,26;2)$, $(12,30;2)$, $(13,18;2)$, $(13,21;2)$, $(13,22;2)$, $(13,25;1)$, $(13,26;2)$, $(14,15;2)$, $(14,21;2)$, $(14,22;1)$, $(14,26;1)$, $(15,22;2)$, $(15,27;1)$, $(16,21;2)$, $(16,25;1)$, $(16,26;1)$, $(16,27;2)$, $(16,28;2)$, $(16,29;2)$, $(17,19;2)$, $(17,20;2)$, $(17,21;1)$, $(17,22;2)$, $(17,30;2)$, $(18,19;2)$, $(18,23;2)$, $(18,24;1)$, $(18,26;1)$, $(18,30;2)$, $(19,21;2)$, $(20,21;2)$, $(21,26;2)$, $(21,31;1)$, $(22,28;2)$, $(23,26;2)$, $(23,28;2)$, $(24,28;1)$, $(24,31;2)$, $(25,30;2)$, $(26,28;1)$, $(26,30;1)$, $(26,31;2)$, $(27,28;1)$.
Let $(\xi,\eta)$ denote $(x,y)$ or $(y,x)$. We call a vertex $i$ \emph{linear} if $H_i=\varepsilon_i(\eta-r_i(\xi))$ for a unit $\varepsilon_i$ of $\Z[\omega]$ and a polynomial $r_i$ with coefficients in $\Z[\omega]$; the last column of Table~\ref{tab:vert1} and its analogues gives $\eta=r_i(\xi)$ for the linear vertices. For a linear vertex $i$, we put $f_i(t)=F$ and $g_{ij}(t)=H_j$ evaluated at $\xi=t$, $\eta=r_i(t)$; these are polynomials in $t$ with coefficients in $\Z[\omega]$.
For every linear vertex $i$, Table~\ref{tab:star} gives an integer $D_i\geq1$ and a polynomial $w_i(t)$ with coefficients in $\Z[\omega]$ such that
\begin{equation}\label{eq:starlin}
D_i^3\,c_i\prod_jg_{ij}^{\rho_{ij}}-w_i^3=f_iQ_i
\end{equation}
for a polynomial $Q_i(t)$ with coefficients in $\Z[\omega]$, where $\rho_{ij}\in\{0,1,2\}$; $Q_i$ is the quotient of the division of the left side by $f_i$, and the remainder is zero.
For every edge $(i,j)$ with a linear endpoint, Table~\ref{tab:sep} gives the first linear endpoint $k\in\{i,j\}$ and the factorization of the norm of the resultant $R_{ij}=\operatorname{Res}_t(f_k,g_{kl})$, where $\{k,l\}=\{i,j\}$; it is non-zero.
These identities can be checked by expanding both sides and by polynomial division, using $\omega^2=-1-\omega$. They hold for each certificate, with its own data.

\section{Proof of Theorem~\ref{thm:main}}

Suppose that $(x,y)\in\Z^2$ is a solution of \eqref{eq:main}, so that $F(x,y)=0$; from now on the polynomials denote their values at $(x,y)$, which lie in $\Z[\omega]$.
\paragraph{Step 1: the values are non-zero.} The vertex $H_{11}$ of the first certificate is linear, and $H_{11}=0$ would give $f_{11}(\xi)=0$ with $\xi\in\Z$, where $f_{11}(t)=\langle 2t^{4}+t^{3}+4t^{2}+4t+1,0\rangle$. An integer root of $2t^{4}+t^{3}+4t^{2}+4t+1$ divides its constant term $1$, and $f_{11}(c)\neq0$ for $c\in\{-1,1\}$; hence $H_{11}\neq0$. For $i\in\{0,\allowbreak 1,\allowbreak 2,\allowbreak 3,\allowbreak 4,\allowbreak 5,\allowbreak 7,\allowbreak 8,\allowbreak 9,\allowbreak 12,\allowbreak 13,\allowbreak 14,\allowbreak 15,\allowbreak 16,\allowbreak 17,\allowbreak 19,\allowbreak 20,\allowbreak 21,\allowbreak 23,\allowbreak 24,\allowbreak 25,\allowbreak 26,\allowbreak 27\}$ in the first certificate, there is no residue pair $(a,b)$ modulo $2$ with $F(a,b)\equiv0$ and $H_i(a,b)\equiv0\pmod{2}$, as one checks for the $2^2$ residue pairs; hence $H_i\neq0$. For $i\in\{10,\allowbreak 28\}$ in the first certificate, there is no residue pair $(a,b)$ modulo $3$ with $F(a,b)\equiv0$ and $H_i(a,b)\equiv0\pmod{3}$, as one checks for the $3^2$ residue pairs; hence $H_i\neq0$. For $i\in\{18,\allowbreak 22\}$ in the first certificate, there is no residue pair $(a,b)$ modulo $4$ with $F(a,b)\equiv0$ and $H_i(a,b)\equiv0\pmod{4}$, as one checks for the $4^2$ residue pairs; hence $H_i\neq0$. For $i\in\{6\}$ in the first certificate, there is no residue pair $(a,b)$ modulo $7$ with $F(a,b)\equiv0$ and $H_i(a,b)\equiv0\pmod{7}$, as one checks for the $7^2$ residue pairs; hence $H_i\neq0$. For $i\in\{1,\allowbreak 2,\allowbreak 5,\allowbreak 6,\allowbreak 7,\allowbreak 9,\allowbreak 10,\allowbreak 12,\allowbreak 14,\allowbreak 15,\allowbreak 16,\allowbreak 17,\allowbreak 20,\allowbreak 22,\allowbreak 23,\allowbreak 24,\allowbreak 27,\allowbreak 28,\allowbreak 29,\allowbreak 30\}$ in the second certificate, there is no residue pair $(a,b)$ modulo $2$ with $F(a,b)\equiv0$ and $H_i(a,b)\equiv0\pmod{2}$, as one checks for the $2^2$ residue pairs; hence $H_i\neq0$. For $i\in\{0,\allowbreak 8,\allowbreak 18,\allowbreak 19,\allowbreak 21,\allowbreak 25,\allowbreak 31\}$ in the second certificate, there is no residue pair $(a,b)$ modulo $3$ with $F(a,b)\equiv0$ and $H_i(a,b)\equiv0\pmod{3}$, as one checks for the $3^2$ residue pairs; hence $H_i\neq0$. For $i\in\{3,\allowbreak 4,\allowbreak 11\}$ in the second certificate, there is no residue pair $(a,b)$ modulo $4$ with $F(a,b)\equiv0$ and $H_i(a,b)\equiv0\pmod{4}$, as one checks for the $4^2$ residue pairs; hence $H_i\neq0$. For $i\in\{26\}$ in the second certificate, there is no residue pair $(a,b)$ modulo $5$ with $F(a,b)\equiv0$ and $H_i(a,b)\equiv0\pmod{5}$, as one checks for the $5^2$ residue pairs; hence $H_i\neq0$. For $i\in\{13\}$ in the second certificate, there is no residue pair $(a,b)$ modulo $7$ with $F(a,b)\equiv0$ and $H_i(a,b)\equiv0\pmod{7}$, as one checks for the $7^2$ residue pairs; hence $H_i\neq0$. Also, all $c_i$ are non-zero.

\paragraph{Step 2: reciprocity.} For every place $v$ of $K$, put
\begin{equation}\label{eq:B}
B^{(g)}_v=\sum_{(i,j)}w_{ij}\,\beta_v(H_i,H_j)+\sum_i\beta_v(H_i,c_i)\in\Z/3\Z,
\end{equation}
the first sum over the edges of the $g$-th certificate, and, for a rational prime $p$, put $B^{(g)}_p=\sum_{v\mid p}B^{(g)}_v$. Applying (H3) to every pair $(H_i,H_j)$ with an edge, multiplied by $w_{ij}$, and to every pair $(H_i,c_i)$, and adding, we find that $B^{(g)}_v=0$ for all but finitely many $v$ and $\sum_pB^{(g)}_p=0$ for every $g$, as the complex place contributes $0$. By (H1) and the definition of $c_p$, $B^{(g)}_p=\sum_{(i,j)}w_{ij}\beta_p(c_p(H_i),c_p(H_j))+\sum_i\beta_p(c_p(H_i),c_p(c_i))$, where $\beta_p$ is the form of Lemma~\ref{lem:M3} if $p=3$.

\paragraph{Step 3: the primes outside $S$.} Let $S=\{2,3,5,7,13\}$. It contains $3$ and the primes dividing the norms of the constants $c_i$, of the $D_i$, of the $R_{ij}$. Let $p\notin S$, and let $v$ be a place above $p$, with residue field $k_v$; then all these numbers are units at $v$. First, no edge $(i,j)$ has both values divisible by $v$. Indeed, suppose that $v$ divides $H_i$ and $H_j$. If $R_{ij}$ is defined, let $k$ be the linear endpoint used for it and $\{k,l\}=\{i,j\}$, and put $t=\xi$. As $H_k=\varepsilon_k(\eta-r_k(\xi))$, we have $\eta\equiv r_k(t)$ modulo $v$, hence $f_k(t)\equiv F(x,y)=0$ and $g_{kl}(t)\equiv H_l\equiv0$ modulo $v$. The resultant can be written as $R_{ij}=Af_k+Cg_{kl}$ with polynomials $A$ and $C$ with coefficients in $\Z[\omega]$, so $v$ divides $R_{ij}$, a contradiction. By (H2), a term of \eqref{eq:B} in which all arguments are units vanishes. Now let $v$ divide $H_i$. Collecting the terms of \eqref{eq:B} which contain $H_i$, and using (H1) and $\rho_{ji}=-w_{ij}$ for the edges $(j,i)$ with $j<i$, we obtain $\beta_v(H_i,\Pi_i)$ with $\Pi_i=c_i\prod_jH_j^{\rho_{ij}}$, which is a unit. We claim that $D_i^3\Pi_i\equiv W^3$ modulo $v$ for some $W\in\Z[\omega]$. If $i$ is linear, put $t=\xi$; then $g_{ij}(t)\equiv H_j$ and $f_i(t)\equiv0$ modulo $v$ as above, and \eqref{eq:starlin} gives the claim with $W=w_i(t)$. Hence the reduction of $\Pi_i$ is a non-zero cube in $k_v$, $\chi(\Pi_i)=0$, and $\beta_v(H_i,\Pi_i)=v(H_i)\chi(\Pi_i)=0$ by (H2). Hence $B^{(g)}_v=0$, and $B^{(g)}_p=0$.

\paragraph{Step 4: the primes in $S$.} Let $p\in S$. Here $2$ is inert, $5$ is inert, the places above $7$ are $v_{2}$ and $v_{4}$, and the places above $13$ are $v_{3}$ and $v_{9}$. Suppose that $x\equiv a$ and $y\equiv b\pmod{p^k}$. Then $H_i\equiv H_i(a,b)\pmod{p^k}$, and Lemma~\ref{lem:residue}, applied with $\zeta$ the residue of $H_i(a,b)$, gives a vector $b_i$ and a subspace $D_i$ with $c_p(H_i)\in b_i+D_i$; when $H_i$ has rational coefficients, $H_i\in\Z$, and we use this in (i) and (iv). With $\gamma_i=c_p(c_i)$, if the hypotheses of Lemma~\ref{lem:envelope} hold, then $B^{(g)}_p$ is determined by $a$, $b$ and $k$. Starting from the solutions of \eqref{eq:main} modulo $p$, we refine a residue class $x\equiv a$, $y\equiv b\pmod{p^k}$ to the classes modulo $p^{k+1}$ that contain solutions of \eqref{eq:main} modulo $p^{k+1}$, until these hypotheses hold for every certificate, or until no solution remains. Table~\ref{tab:local} lists the resulting terminal classes and the value of the vector $(B^{(g)}_p)_g$ in each; every integer solution lies in one of them. The result is $(B^{(g)}_{2})_g=(1,2)$, $(B^{(g)}_{3})_g\in\{(0,1),\allowbreak (0,2),\allowbreak (1,1),\allowbreak (1,2),\allowbreak (2,0)\}$, $(B^{(g)}_{5})_g=(0,0)$, $(B^{(g)}_{7})_g=(0,0)$, $(B^{(g)}_{13})_g=(0,0)$.

\paragraph{Step 5: contradiction.} By Steps~2 and~3, $\sum_{p\in S}B^{(g)}_p=0$ for every $g$. But by Step~4, the vector $(\sum_{p\in S}B^{(g)}_p)_g$ lies in $\{(0,2),\allowbreak (1,0),\allowbreak (1,1),\allowbreak (2,0),\allowbreak (2,1)\}$, which does not contain $(0,0)$. This contradiction proves the theorem.
\qed


\begin{thebibliography}{9}
\bibitem{N} J.~Neukirch, \emph{Algebraic Number Theory}, Grundlehren der mathematischen Wissenschaften 322, Springer, Berlin, 1999.
\end{thebibliography}

\clearpage
\begingroup\small
\setlength{\LTcapwidth}{\textwidth}
\begin{longtable}{@{}rlll@{}}
\caption{The polynomials $H_i$, the constants $c_i$ and, for the linear vertices, $\eta=r_i(\xi)$ of the first certificate. Here $\langle A,B\rangle=A+B\omega$.}\label{tab:vert1}\\
\hline $i$ & $H_i$ & $c_i$ & chart\\\hline\endfirsthead\hline $i$ & $H_i$ & $c_i$ & chart\\\hline\endhead\hline\endfoot
0 & $\langle x^{2}-y,0\rangle$ & $\langle 1,0\rangle$ & $y=\langle x^{2},0\rangle$\\
1 & $\langle x-y+1,-x^{2}+x+1\rangle$ & $\langle 1,0\rangle$ & $y=\langle x+1,-x^{2}+x+1\rangle$\\
2 & $\langle -x^{2}-y+1,-x^{2}+x+1\rangle$ & $\langle 1,0\rangle$ & $y=\langle -x^{2}+1,-x^{2}+x+1\rangle$\\
3 & $\langle x,0\rangle$ & $\langle 0,-1\rangle$ & $x=\langle 0,0\rangle$\\
4 & $\langle x+1,1\rangle$ & $\langle 0,-1\rangle$ & $x=\langle -1,-1\rangle$\\
5 & $\langle x,-1\rangle$ & $\langle 1,-1\rangle$ & $x=\langle 0,1\rangle$\\
6 & $\langle x-1,0\rangle$ & $\langle -1,-2\rangle$ & $x=\langle 1,0\rangle$\\
7 & $\langle x-1,-1\rangle$ & $\langle 0,-1\rangle$ & $x=\langle 1,1\rangle$\\
8 & $\langle x,-y\rangle$ & $\langle 4,0\rangle$ & $y=\langle -x,-x\rangle$\\
9 & $\langle x+y,0\rangle$ & $\langle 1,0\rangle$ & $y=\langle -x,0\rangle$\\
10 & $\langle x-y+1,-y\rangle$ & $\langle 2,0\rangle$ & $y=\langle 0,-x-1\rangle$\\
11 & $\langle x-y+1,0\rangle$ & $\langle 1,0\rangle$ & $y=\langle x+1,0\rangle$\\
12 & $\langle x+y-1,y-1\rangle$ & $\langle -1,-1\rangle$ & $y=\langle 1,x\rangle$\\
13 & $\langle x^{2}-y,x^{2}\rangle$ & $\langle -1,-1\rangle$ & $y=\langle x^{2},x^{2}\rangle$\\
14 & $\langle x-y,-y+1\rangle$ & $\langle -1,-1\rangle$ & $y=\langle 1,-x+1\rangle$\\
15 & $\langle x+y,y-1\rangle$ & $\langle 4,0\rangle$ & $y=\langle 1,x+1\rangle$\\
16 & $\langle -y+1,1\rangle$ & $\langle -3,-3\rangle$ & $y=\langle 1,1\rangle$\\
17 & $\langle -y+1,0\rangle$ & $\langle -1,-1\rangle$ & $y=\langle 1,0\rangle$\\
18 & $\langle -2x-y,-x-1\rangle$ & $\langle 4,-4\rangle$ & $y=\langle -2x,-x-1\rangle$\\
19 & $\langle -x^{2}-y,-x^{2}\rangle$ & $\langle 1,0\rangle$ & $y=\langle -x^{2},-x^{2}\rangle$\\
20 & $\langle -x^{2}-y+1,-x^{2}\rangle$ & $\langle -1,-1\rangle$ & $y=\langle -x^{2}+1,-x^{2}\rangle$\\
21 & $\langle x^{2}-x-y,x^{2}-x-1\rangle$ & $\langle -1,-1\rangle$ & $y=\langle x^{2}-x,x^{2}-x-1\rangle$\\
22 & $\langle x+1,0\rangle$ & $\langle -2,-2\rangle$ & $x=\langle -1,0\rangle$\\
23 & $\langle x,y\rangle$ & $\langle 0,-2\rangle$ & $y=\langle x,x\rangle$\\
24 & $\langle x-y,-y-1\rangle$ & $\langle 0,-2\rangle$ & $y=\langle -1,-x-1\rangle$\\
25 & $\langle -y,x\rangle$ & $\langle 0,-1\rangle$ & $y=\langle 0,x\rangle$\\
26 & $\langle x,1\rangle$ & $\langle -1,-1\rangle$ & $x=\langle 0,-1\rangle$\\
27 & $\langle -x-y,x^{2}-x-1\rangle$ & $\langle 1,0\rangle$ & $y=\langle -x,x^{2}-x-1\rangle$\\
28 & $\langle 2x-y,x+1\rangle$ & $\langle 1,0\rangle$ & $y=\langle 2x,x+1\rangle$\\
\end{longtable}
\endgroup

\begingroup\small
\setlength{\LTcapwidth}{\textwidth}
\begin{longtable}{@{}rlll@{}}
\caption{The polynomials $H_i$, the constants $c_i$ and, for the linear vertices, $\eta=r_i(\xi)$ of the second certificate. Here $\langle A,B\rangle=A+B\omega$.}\label{tab:vert2}\\
\hline $i$ & $H_i$ & $c_i$ & chart\\\hline\endfirsthead\hline $i$ & $H_i$ & $c_i$ & chart\\\hline\endhead\hline\endfoot
0 & $\langle -y,0\rangle$ & $\langle 4,0\rangle$ & $y=\langle 0,0\rangle$\\
1 & $\langle -y+1,1\rangle$ & $\langle -4,-8\rangle$ & $y=\langle 1,1\rangle$\\
2 & $\langle -y,-1\rangle$ & $\langle 4,0\rangle$ & $y=\langle 0,-1\rangle$\\
3 & $\langle -x-y+1,x+1\rangle$ & $\langle -4,-2\rangle$ & $y=\langle -x+1,x+1\rangle$\\
4 & $\langle -2x-y,-x-1\rangle$ & $\langle 2,-2\rangle$ & $y=\langle -2x,-x-1\rangle$\\
5 & $\langle x^{2}-y,0\rangle$ & $\langle 4,0\rangle$ & $y=\langle x^{2},0\rangle$\\
6 & $\langle -x^{2}-y,-x^{2}\rangle$ & $\langle 1,0\rangle$ & $y=\langle -x^{2},-x^{2}\rangle$\\
7 & $\langle x^{2}-y,x^{2}+x\rangle$ & $\langle 1,0\rangle$ & $y=\langle x^{2},x^{2}+x\rangle$\\
8 & $\langle -x^{2}-y+1,-x^{2}-x\rangle$ & $\langle 1,0\rangle$ & $y=\langle -x^{2}+1,-x^{2}-x\rangle$\\
9 & $\langle x,0\rangle$ & $\langle 1,0\rangle$ & $x=\langle 0,0\rangle$\\
10 & $\langle x+1,1\rangle$ & $\langle 0,-3\rangle$ & $x=\langle -1,-1\rangle$\\
11 & $\langle x+1,0\rangle$ & $\langle 4,0\rangle$ & $x=\langle -1,0\rangle$\\
12 & $\langle x,-1\rangle$ & $\langle -1,-2\rangle$ & $x=\langle 0,1\rangle$\\
13 & $\langle x-1,0\rangle$ & $\langle 0,-3\rangle$ & $x=\langle 1,0\rangle$\\
14 & $\langle x-1,-1\rangle$ & $\langle 1,0\rangle$ & $x=\langle 1,1\rangle$\\
15 & $\langle x-y,-y\rangle$ & $\langle 1,0\rangle$ & $y=\langle 0,-x\rangle$\\
16 & $\langle x,-y\rangle$ & $\langle 1,0\rangle$ & $y=\langle -x,-x\rangle$\\
17 & $\langle x+y,0\rangle$ & $\langle 1,0\rangle$ & $y=\langle -x,0\rangle$\\
18 & $\langle x-y+1,-y\rangle$ & $\langle -1,-1\rangle$ & $y=\langle 0,-x-1\rangle$\\
19 & $\langle x+y+1,y\rangle$ & $\langle 0,-1\rangle$ & $y=\langle 0,x+1\rangle$\\
20 & $\langle x-y,-y-1\rangle$ & $\langle -2,-2\rangle$ & $y=\langle -1,-x-1\rangle$\\
21 & $\langle x-y-1,0\rangle$ & $\langle 1,0\rangle$ & $y=\langle x-1,0\rangle$\\
22 & $\langle x+y-1,y-1\rangle$ & $\langle 2,0\rangle$ & $y=\langle 1,x\rangle$\\
23 & $\langle -y,x\rangle$ & $\langle 1,0\rangle$ & $y=\langle 0,x\rangle$\\
24 & $\langle -y,-x^{2}\rangle$ & $\langle 1,0\rangle$ & $y=\langle 0,-x^{2}\rangle$\\
25 & $\langle x-y+1,x+1\rangle$ & $\langle 0,-2\rangle$ & $y=\langle x+1,x+1\rangle$\\
26 & $\langle x^{2}+x-y,0\rangle$ & $\langle 0,-1\rangle$ & $y=\langle x^{2}+x,0\rangle$\\
27 & $\langle -y+1,x+1\rangle$ & $\langle -4,-4\rangle$ & $y=\langle 1,x+1\rangle$\\
28 & $\langle x-y+1,-y+1\rangle$ & $\langle 4,-4\rangle$ & $y=\langle 1,-x\rangle$\\
29 & $\langle x,1\rangle$ & $\langle 1,0\rangle$ & $x=\langle 0,-1\rangle$\\
30 & $\langle -x^{2}-y,0\rangle$ & $\langle 4,0\rangle$ & $y=\langle -x^{2},0\rangle$\\
31 & $\langle x-y-1,-y\rangle$ & $\langle 0,-6\rangle$ & $y=\langle 0,-x+1\rangle$\\
\end{longtable}
\endgroup

\begingroup\scriptsize
\setlength{\LTcapwidth}{\textwidth}
\begin{longtable}{@{}ll@{}}
\caption{The data of the identities \eqref{eq:starlin}; the entry $j{:}e$ means $\rho_{ij}=e\neq0$.}\label{tab:star}\\
\hline\endfirsthead\hline\endhead\hline\endfoot
\multicolumn{2}{@{}l}{\emph{The first certificate.}}\\
\multicolumn{2}{@{}p{0.97\textwidth}@{}}{$i=0$, $D_i=4$, $\rho_{ij}$: $1{:}2, \allowbreak 4{:}1, \allowbreak 5{:}1, \allowbreak 7{:}1, \allowbreak 10{:}1, \allowbreak 13{:}2, \allowbreak 15{:}1, \allowbreak 16{:}2, \allowbreak 18{:}1, \allowbreak 27{:}2$}\\
$w_i$ & $\langle -4t^{5}-36t^{4}-4t^{3}+20t^{2}-12t-12,$\\
 & $\quad 12t^{5}+40t^{4}-12t^{2}+24t+16\rangle$\\
\addlinespace[2pt]
\multicolumn{2}{@{}p{0.97\textwidth}@{}}{$i=1$, $D_i=1$, $\rho_{ij}$: $0{:}1, \allowbreak 2{:}2, \allowbreak 3{:}2, \allowbreak 5{:}1, \allowbreak 8{:}1, \allowbreak 11{:}1, \allowbreak 12{:}2, \allowbreak 15{:}1, \allowbreak 23{:}1$}\\
$w_i$ & $\langle -8t^{5}-7t^{4}+6t^{3}+12t^{2}+5t,$\\
 & $\quad -4t^{5}-13t^{4}-11t^{3}+7t+3\rangle$\\
\addlinespace[2pt]
\multicolumn{2}{@{}p{0.97\textwidth}@{}}{$i=2$, $D_i=18$, $\rho_{ij}$: $1{:}1, \allowbreak 5{:}2, \allowbreak 9{:}2, \allowbreak 16{:}2, \allowbreak 17{:}2$}\\
$w_i$ & $\langle -18t^{4}-18t^{3},$\\
 & $\quad -18t^{4}-18t^{3}\rangle$\\
\addlinespace[2pt]
\multicolumn{2}{@{}p{0.97\textwidth}@{}}{$i=3$, $D_i=2$, $\rho_{ij}$: $1{:}1, \allowbreak 12{:}2, \allowbreak 16{:}1, \allowbreak 18{:}2, \allowbreak 19{:}1, \allowbreak 23{:}1$}\\
$w_i$ & $\langle 2t^{2}-4t+2,$\\
 & $\quad 0\rangle$\\
\addlinespace[2pt]
\multicolumn{2}{@{}p{0.97\textwidth}@{}}{$i=4$, $D_i=2$, $\rho_{ij}$: $0{:}2, \allowbreak 18{:}2, \allowbreak 23{:}1, \allowbreak 24{:}1, \allowbreak 27{:}2$}\\
$w_i$ & $\langle -2t^{2}-4t+2,$\\
 & $\quad -4t-2\rangle$\\
\addlinespace[2pt]
\multicolumn{2}{@{}p{0.97\textwidth}@{}}{$i=5$, $D_i=4$, $\rho_{ij}$: $0{:}2, \allowbreak 1{:}2, \allowbreak 2{:}1, \allowbreak 9{:}1, \allowbreak 19{:}1, \allowbreak 23{:}2$}\\
$w_i$ & $\langle 4t^{2}-4t-12,$\\
 & $\quad -8t^{2}-16t-16\rangle$\\
\addlinespace[2pt]
\multicolumn{2}{@{}p{0.97\textwidth}@{}}{$i=6$, $D_i=2$, $\rho_{ij}$: $9{:}2, \allowbreak 10{:}1, \allowbreak 11{:}2, \allowbreak 18{:}1, \allowbreak 21{:}1, \allowbreak 28{:}2$}\\
$w_i$ & $\langle 2t^{2}-4t,$\\
 & $\quad -2t^{2}+4t\rangle$\\
\addlinespace[2pt]
\multicolumn{2}{@{}p{0.97\textwidth}@{}}{$i=7$, $D_i=14$, $\rho_{ij}$: $0{:}2, \allowbreak 9{:}2, \allowbreak 14{:}1, \allowbreak 19{:}1$}\\
$w_i$ & $\langle -10t^{2}+20t+10,$\\
 & $\quad -8t^{2}+16t-6\rangle$\\
\addlinespace[2pt]
\multicolumn{2}{@{}p{0.97\textwidth}@{}}{$i=8$, $D_i=1$, $\rho_{ij}$: $1{:}2, \allowbreak 9{:}1, \allowbreak 15{:}2, \allowbreak 16{:}1$}\\
$w_i$ & $\langle -2t^{2}+2t,$\\
 & $\quad 2t^{3}-2t-2\rangle$\\
\addlinespace[2pt]
\multicolumn{2}{@{}p{0.97\textwidth}@{}}{$i=9$, $D_i=8$, $\rho_{ij}$: $2{:}1, \allowbreak 5{:}2, \allowbreak 6{:}1, \allowbreak 7{:}1, \allowbreak 8{:}2, \allowbreak 14{:}2, \allowbreak 20{:}2, \allowbreak 22{:}1$}\\
$w_i$ & $\langle -4t^{3}-12t^{2}-16t-4,$\\
 & $\quad 2t^{3}-10t^{2}-16t-6\rangle$\\
\addlinespace[2pt]
\multicolumn{2}{@{}p{0.97\textwidth}@{}}{$i=10$, $D_i=2$, $\rho_{ij}$: $0{:}2, \allowbreak 6{:}2, \allowbreak 20{:}1, \allowbreak 22{:}1, \allowbreak 25{:}2$}\\
$w_i$ & $\langle -4t^{3}-4t^{2}+4t+4,$\\
 & $\quad -4t^{3}+8t+4\rangle$\\
\addlinespace[2pt]
\multicolumn{2}{@{}p{0.97\textwidth}@{}}{$i=11$, $D_i=4$, $\rho_{ij}$: $1{:}2, \allowbreak 6{:}1, \allowbreak 15{:}2, \allowbreak 16{:}1, \allowbreak 17{:}2, \allowbreak 22{:}1, \allowbreak 27{:}2$}\\
$w_i$ & $\langle 15t^{3}+8t^{2}-2t-1,$\\
 & $\quad 17t^{3}+24t^{2}+10t+1\rangle$\\
\addlinespace[2pt]
\multicolumn{2}{@{}p{0.97\textwidth}@{}}{$i=12$, $D_i=2$, $\rho_{ij}$: $1{:}1, \allowbreak 3{:}1, \allowbreak 13{:}1, \allowbreak 14{:}1, \allowbreak 28{:}2$}\\
$w_i$ & $\langle -2t+4,$\\
 & $\quad -4t^{2}+2\rangle$\\
\addlinespace[2pt]
\multicolumn{2}{@{}p{0.97\textwidth}@{}}{$i=13$, $D_i=4$, $\rho_{ij}$: $0{:}1, \allowbreak 12{:}2, \allowbreak 14{:}1, \allowbreak 15{:}2, \allowbreak 16{:}2, \allowbreak 21{:}2, \allowbreak 27{:}2$}\\
$w_i$ & $\langle 4t^{5}+4t^{3}+8t^{2}-4,$\\
 & $\quad 4t^{5}-4t^{3}+8t^{2}+4t\rangle$\\
\addlinespace[2pt]
\multicolumn{2}{@{}p{0.97\textwidth}@{}}{$i=14$, $D_i=2$, $\rho_{ij}$: $7{:}2, \allowbreak 9{:}1, \allowbreak 12{:}2, \allowbreak 13{:}2, \allowbreak 15{:}1$}\\
$w_i$ & $\langle 4t^{2}-4t,$\\
 & $\quad 8t^{2}\rangle$\\
\addlinespace[2pt]
\multicolumn{2}{@{}p{0.97\textwidth}@{}}{$i=15$, $D_i=2$, $\rho_{ij}$: $0{:}2, \allowbreak 1{:}2, \allowbreak 8{:}1, \allowbreak 11{:}1, \allowbreak 13{:}1, \allowbreak 14{:}2, \allowbreak 19{:}1$}\\
$w_i$ & $\langle -5t^{3}-6t^{2}-t+1,$\\
 & $\quad -8t^{3}-8t^{2}+t+1\rangle$\\
\addlinespace[2pt]
\multicolumn{2}{@{}p{0.97\textwidth}@{}}{$i=16$, $D_i=2$, $\rho_{ij}$: $0{:}1, \allowbreak 2{:}1, \allowbreak 3{:}2, \allowbreak 8{:}2, \allowbreak 11{:}2, \allowbreak 13{:}1$}\\
$w_i$ & $\langle 6t^{3}+2t^{2}-3t+1,$\\
 & $\quad -2t^{2}-3t+2\rangle$\\
\addlinespace[2pt]
\multicolumn{2}{@{}p{0.97\textwidth}@{}}{$i=17$, $D_i=2$, $\rho_{ij}$: $2{:}1, \allowbreak 11{:}1, \allowbreak 21{:}2, \allowbreak 27{:}1$}\\
$w_i$ & $\langle 2t^{3}+2t^{2},$\\
 & $\quad 0\rangle$\\
\addlinespace[2pt]
\multicolumn{2}{@{}p{0.97\textwidth}@{}}{$i=18$, $D_i=2$, $\rho_{ij}$: $0{:}2, \allowbreak 3{:}1, \allowbreak 4{:}1, \allowbreak 6{:}2, \allowbreak 24{:}2, \allowbreak 26{:}1$}\\
$w_i$ & $\langle 8t^{3}+24t^{2}-8t,$\\
 & $\quad -8t^{3}+24t^{2}+8t\rangle$\\
\addlinespace[2pt]
\multicolumn{2}{@{}p{0.97\textwidth}@{}}{$i=19$, $D_i=2$, $\rho_{ij}$: $3{:}2, \allowbreak 5{:}2, \allowbreak 7{:}2, \allowbreak 15{:}2, \allowbreak 24{:}1, \allowbreak 26{:}2$}\\
$w_i$ & $\langle -2t^{2}+2,$\\
 & $\quad -2t^{2}+2t+2\rangle$\\
\addlinespace[2pt]
\multicolumn{2}{@{}p{0.97\textwidth}@{}}{$i=20$, $D_i=26$, $\rho_{ij}$: $9{:}1, \allowbreak 10{:}2, \allowbreak 21{:}1$}\\
$w_i$ & $\langle -8t^{5}+24t^{4}-30t^{3}-8t^{2}+62t+32,$\\
 & $\quad 24t^{5}+32t^{4}-40t^{3}+24t^{2}+48t+8\rangle$\\
\addlinespace[2pt]
\multicolumn{2}{@{}p{0.97\textwidth}@{}}{$i=21$, $D_i=42$, $\rho_{ij}$: $6{:}2, \allowbreak 13{:}1, \allowbreak 17{:}1, \allowbreak 20{:}2, \allowbreak 25{:}1$}\\
$w_i$ & $\langle -84t^{3}+42t^{2}+42t,$\\
 & $\quad 0\rangle$\\
\addlinespace[2pt]
\multicolumn{2}{@{}p{0.97\textwidth}@{}}{$i=22$, $D_i=2$, $\rho_{ij}$: $9{:}2, \allowbreak 10{:}2, \allowbreak 11{:}2, \allowbreak 28{:}2$}\\
$w_i$ & $\langle 4,$\\
 & $\quad 0\rangle$\\
\addlinespace[2pt]
\multicolumn{2}{@{}p{0.97\textwidth}@{}}{$i=23$, $D_i=2$, $\rho_{ij}$: $1{:}2, \allowbreak 3{:}2, \allowbreak 4{:}2, \allowbreak 5{:}1, \allowbreak 26{:}1$}\\
$w_i$ & $\langle -2t^{3}-2t^{2}+2,$\\
 & $\quad -2t^{2}-2t\rangle$\\
\addlinespace[2pt]
\multicolumn{2}{@{}p{0.97\textwidth}@{}}{$i=24$, $D_i=6$, $\rho_{ij}$: $4{:}2, \allowbreak 18{:}1, \allowbreak 19{:}2, \allowbreak 26{:}1, \allowbreak 27{:}1$}\\
$w_i$ & $\langle 4t^{2}+4,$\\
 & $\quad -4t^{3}-4t^{2}-16t\rangle$\\
\addlinespace[2pt]
\multicolumn{2}{@{}p{0.97\textwidth}@{}}{$i=25$, $D_i=52$, $\rho_{ij}$: $10{:}1, \allowbreak 21{:}2$}\\
$w_i$ & $\langle 28t^{3}-30t^{2}+32t+16,$\\
 & $\quad 60t^{3}+62t^{2}+76t+38\rangle$\\
\addlinespace[2pt]
\multicolumn{2}{@{}p{0.97\textwidth}@{}}{$i=26$, $D_i=14$, $\rho_{ij}$: $18{:}2, \allowbreak 19{:}1, \allowbreak 23{:}2, \allowbreak 24{:}2, \allowbreak 28{:}2$}\\
$w_i$ & $\langle -10t^{2}-8t+24,$\\
 & $\quad -16t^{2}-10t+2\rangle$\\
\addlinespace[2pt]
\multicolumn{2}{@{}p{0.97\textwidth}@{}}{$i=27$, $D_i=18$, $\rho_{ij}$: $0{:}1, \allowbreak 4{:}1, \allowbreak 11{:}1, \allowbreak 13{:}1, \allowbreak 17{:}2, \allowbreak 24{:}2$}\\
$w_i$ & $\langle -36t^{4}-36t^{3}+36t^{2}+54t+18,$\\
 & $\quad -18t^{4}-72t^{3}-54t^{2}+18t+18\rangle$\\
\addlinespace[2pt]
\multicolumn{2}{@{}p{0.97\textwidth}@{}}{$i=28$, $D_i=2$, $\rho_{ij}$: $6{:}1, \allowbreak 12{:}1, \allowbreak 22{:}1, \allowbreak 26{:}1$}\\
$w_i$ & $\langle 8t^{2}-2t-2,$\\
 & $\quad -4t^{3}+6t^{2}+4t-2\rangle$\\
\addlinespace[2pt]
\multicolumn{2}{@{}l}{\emph{The second certificate.}}\\
\multicolumn{2}{@{}p{0.97\textwidth}@{}}{$i=0$, $D_i=1$, $\rho_{ij}$: $15{:}1, \allowbreak 26{:}1$}\\
$w_i$ & $\langle 2t^{2},$\\
 & $\quad 2t^{2}\rangle$\\
\addlinespace[2pt]
\multicolumn{2}{@{}p{0.97\textwidth}@{}}{$i=1$, $D_i=1$, $\rho_{ij}$: $6{:}2, \allowbreak 11{:}2, \allowbreak 12{:}1, \allowbreak 22{:}2$}\\
$w_i$ & $\langle 2t+1,$\\
 & $\quad t-1\rangle$\\
\addlinespace[2pt]
\multicolumn{2}{@{}p{0.97\textwidth}@{}}{$i=2$, $D_i=4$, $\rho_{ij}$: $5{:}2, \allowbreak 6{:}1, \allowbreak 7{:}2, \allowbreak 8{:}1, \allowbreak 10{:}1, \allowbreak 11{:}1, \allowbreak 13{:}1, \allowbreak 17{:}1, \allowbreak 18{:}1, \allowbreak 19{:}2, \allowbreak 21{:}1, \allowbreak 26{:}1, \allowbreak 29{:}2$}\\
$w_i$ & $\langle -26t^{3}-70t^{2}-36t+36,$\\
 & $\quad 44t^{3}-26t^{2}-60t+12\rangle$\\
\addlinespace[2pt]
\multicolumn{2}{@{}p{0.97\textwidth}@{}}{$i=3$, $D_i=2$, $\rho_{ij}$: $4{:}2, \allowbreak 8{:}1, \allowbreak 16{:}1, \allowbreak 21{:}2, \allowbreak 23{:}1, \allowbreak 25{:}2$}\\
$w_i$ & $\langle 16t^{3}+16t^{2},$\\
 & $\quad 20t^{3}+8t^{2}-12t\rangle$\\
\addlinespace[2pt]
\multicolumn{2}{@{}p{0.97\textwidth}@{}}{$i=4$, $D_i=2$, $\rho_{ij}$: $3{:}1, \allowbreak 5{:}1, \allowbreak 17{:}1, \allowbreak 20{:}1, \allowbreak 21{:}1, \allowbreak 22{:}1, \allowbreak 23{:}1, \allowbreak 29{:}2$}\\
$w_i$ & $\langle -8t^{3}-32t^{2}+8t+8,$\\
 & $\quad 8t^{3}-28t^{2}-20t+4\rangle$\\
\addlinespace[2pt]
\multicolumn{2}{@{}p{0.97\textwidth}@{}}{$i=5$, $D_i=2$, $\rho_{ij}$: $2{:}1, \allowbreak 4{:}2, \allowbreak 10{:}1, \allowbreak 14{:}2, \allowbreak 16{:}1, \allowbreak 18{:}1, \allowbreak 19{:}1, \allowbreak 23{:}2, \allowbreak 24{:}1, \allowbreak 25{:}1, \allowbreak 31{:}1$}\\
$w_i$ & $\langle -8t^{5}-16t^{4}+8t^{3}+4t^{2}-12t-4,$\\
 & $\quad -4t^{5}-4t^{4}-4t^{2}-4t\rangle$\\
\addlinespace[2pt]
\multicolumn{2}{@{}p{0.97\textwidth}@{}}{$i=6$, $D_i=4$, $\rho_{ij}$: $1{:}1, \allowbreak 2{:}2, \allowbreak 12{:}1, \allowbreak 17{:}2, \allowbreak 19{:}1, \allowbreak 20{:}1, \allowbreak 27{:}2, \allowbreak 30{:}2$}\\
$w_i$ & $\langle -4t^{5}+4t^{4}-4t^{3}-4t^{2}+4,$\\
 & $\quad -4t^{5}-4t^{4}-8t^{3}-4t^{2}-4t\rangle$\\
\addlinespace[2pt]
\multicolumn{2}{@{}p{0.97\textwidth}@{}}{$i=7$, $D_i=2$, $\rho_{ij}$: $2{:}1, \allowbreak 8{:}1, \allowbreak 9{:}2, \allowbreak 14{:}1, \allowbreak 16{:}2, \allowbreak 17{:}1, \allowbreak 19{:}1, \allowbreak 22{:}2, \allowbreak 25{:}1, \allowbreak 29{:}1$}\\
$w_i$ & $\langle -4t^{5}+4t^{4}-2t^{2}+4t+2,$\\
 & $\quad -8t^{5}+2t^{3}-6t^{2}+2\rangle$\\
\addlinespace[2pt]
\multicolumn{2}{@{}p{0.97\textwidth}@{}}{$i=8$, $D_i=12500$, $\rho_{ij}$: $2{:}2, \allowbreak 3{:}2, \allowbreak 7{:}2, \allowbreak 9{:}1, \allowbreak 13{:}1, \allowbreak 15{:}2, \allowbreak 17{:}2, \allowbreak 19{:}2, \allowbreak 21{:}2, \allowbreak 22{:}2, \allowbreak 29{:}1$}\\
$w_i$ & $\langle 425000t^{5}+2012500t^{4}-1612500t^{3}-2175000t^{2}+300000t+300000,$\\
 & $\quad -112500t^{5}+1887500t^{4}+1137500t^{3}-1537500t^{2}-900000t-75000\rangle$\\
\addlinespace[2pt]
\multicolumn{2}{@{}p{0.97\textwidth}@{}}{$i=9$, $D_i=2$, $\rho_{ij}$: $7{:}1, \allowbreak 8{:}2, \allowbreak 21{:}1$}\\
$w_i$ & $\langle -2t,$\\
 & $\quad 0\rangle$\\
\addlinespace[2pt]
\multicolumn{2}{@{}p{0.97\textwidth}@{}}{$i=10$, $D_i=2$, $\rho_{ij}$: $2{:}2, \allowbreak 5{:}2, \allowbreak 15{:}2, \allowbreak 27{:}1, \allowbreak 30{:}2$}\\
$w_i$ & $\langle 2t^{2}+4,$\\
 & $\quad -2t^{2}+2\rangle$\\
\addlinespace[2pt]
\multicolumn{2}{@{}p{0.97\textwidth}@{}}{$i=11$, $D_i=1$, $\rho_{ij}$: $1{:}1, \allowbreak 2{:}2, \allowbreak 18{:}1, \allowbreak 20{:}2, \allowbreak 21{:}2, \allowbreak 24{:}2, \allowbreak 27{:}1$}\\
$w_i$ & $\langle -2t^{2}-2t,$\\
 & $\quad -2t-2\rangle$\\
\addlinespace[2pt]
\multicolumn{2}{@{}p{0.97\textwidth}@{}}{$i=12$, $D_i=4$, $\rho_{ij}$: $1{:}2, \allowbreak 6{:}2, \allowbreak 23{:}1, \allowbreak 25{:}2, \allowbreak 26{:}2, \allowbreak 30{:}2$}\\
$w_i$ & $\langle 12t^{2}+24t+28,$\\
 & $\quad 12t^{2}+24t+8\rangle$\\
\addlinespace[2pt]
\multicolumn{2}{@{}p{0.97\textwidth}@{}}{$i=13$, $D_i=4$, $\rho_{ij}$: $2{:}2, \allowbreak 8{:}2, \allowbreak 18{:}2, \allowbreak 21{:}2, \allowbreak 22{:}2, \allowbreak 25{:}1, \allowbreak 26{:}2$}\\
$w_i$ & $\langle 48t^{2}-96t+144,$\\
 & $\quad 24t^{2}-48t+48\rangle$\\
\addlinespace[2pt]
\multicolumn{2}{@{}p{0.97\textwidth}@{}}{$i=14$, $D_i=2$, $\rho_{ij}$: $5{:}1, \allowbreak 7{:}2, \allowbreak 15{:}2, \allowbreak 21{:}2, \allowbreak 22{:}1, \allowbreak 26{:}1$}\\
$w_i$ & $\langle 4t+2,$\\
 & $\quad 2t^{2}+4t-4\rangle$\\
\addlinespace[2pt]
\multicolumn{2}{@{}p{0.97\textwidth}@{}}{$i=15$, $D_i=2$, $\rho_{ij}$: $0{:}2, \allowbreak 8{:}1, \allowbreak 10{:}1, \allowbreak 14{:}1, \allowbreak 22{:}2, \allowbreak 27{:}1$}\\
$w_i$ & $\langle -4t^{3}+2t^{2},$\\
 & $\quad 2t^{2}\rangle$\\
\addlinespace[2pt]
\multicolumn{2}{@{}p{0.97\textwidth}@{}}{$i=16$, $D_i=4$, $\rho_{ij}$: $3{:}2, \allowbreak 5{:}2, \allowbreak 7{:}1, \allowbreak 21{:}2, \allowbreak 25{:}1, \allowbreak 26{:}1, \allowbreak 27{:}2, \allowbreak 28{:}2, \allowbreak 29{:}2$}\\
$w_i$ & $\langle -29t^{3}-51t^{2}+2t+1,$\\
 & $\quad 12t^{3}-31t^{2}-33t-18\rangle$\\
\addlinespace[2pt]
\multicolumn{2}{@{}p{0.97\textwidth}@{}}{$i=17$, $D_i=160$, $\rho_{ij}$: $2{:}2, \allowbreak 4{:}2, \allowbreak 6{:}1, \allowbreak 7{:}2, \allowbreak 8{:}1, \allowbreak 19{:}2, \allowbreak 20{:}2, \allowbreak 21{:}1, \allowbreak 22{:}2, \allowbreak 30{:}2$}\\
$w_i$ & $\langle 722t^{3}+78t^{2}-960t-398,$\\
 & $\quad 1500t^{3}+420t^{2}-1280t-580\rangle$\\
\addlinespace[2pt]
\multicolumn{2}{@{}p{0.97\textwidth}@{}}{$i=18$, $D_i=4$, $\rho_{ij}$: $2{:}2, \allowbreak 5{:}2, \allowbreak 11{:}2, \allowbreak 13{:}1, \allowbreak 19{:}2, \allowbreak 23{:}2, \allowbreak 24{:}1, \allowbreak 26{:}1, \allowbreak 30{:}2$}\\
$w_i$ & $\langle -39t^{3}-44t^{2}-12t+11,$\\
 & $\quad 17t^{3}+7t^{2}+22t+22\rangle$\\
\addlinespace[2pt]
\multicolumn{2}{@{}p{0.97\textwidth}@{}}{$i=19$, $D_i=8$, $\rho_{ij}$: $2{:}1, \allowbreak 5{:}2, \allowbreak 6{:}2, \allowbreak 7{:}2, \allowbreak 8{:}1, \allowbreak 17{:}1, \allowbreak 18{:}1, \allowbreak 21{:}2$}\\
$w_i$ & $\langle -4t^{3}+57t^{2}+86t+45,$\\
 & $\quad 5t^{3}+40t^{2}+22t+1\rangle$\\
\addlinespace[2pt]
\multicolumn{2}{@{}p{0.97\textwidth}@{}}{$i=20$, $D_i=6$, $\rho_{ij}$: $4{:}2, \allowbreak 6{:}2, \allowbreak 11{:}1, \allowbreak 17{:}1, \allowbreak 21{:}2$}\\
$w_i$ & $\langle -8t^{3}+8t^{2}+28t+16,$\\
 & $\quad 16t^{2}+12t-8\rangle$\\
\addlinespace[2pt]
\multicolumn{2}{@{}p{0.97\textwidth}@{}}{$i=21$, $D_i=8$, $\rho_{ij}$: $2{:}2, \allowbreak 3{:}1, \allowbreak 4{:}2, \allowbreak 8{:}1, \allowbreak 9{:}2, \allowbreak 11{:}1, \allowbreak 13{:}1, \allowbreak 14{:}1, \allowbreak 16{:}1, \allowbreak 17{:}2, \allowbreak 19{:}1, \allowbreak 20{:}1, \allowbreak 26{:}2, \allowbreak 31{:}1$}\\
$w_i$ & $\langle 45t^{3}-42t^{2}-32t+41,$\\
 & $\quad -520t^{3}+552t^{2}-288t-160\rangle$\\
\addlinespace[2pt]
\multicolumn{2}{@{}p{0.97\textwidth}@{}}{$i=22$, $D_i=2$, $\rho_{ij}$: $1{:}1, \allowbreak 4{:}2, \allowbreak 7{:}1, \allowbreak 8{:}1, \allowbreak 13{:}1, \allowbreak 14{:}2, \allowbreak 15{:}1, \allowbreak 17{:}1, \allowbreak 28{:}2$}\\
$w_i$ & $\langle 12t^{3}-20t^{2}+8t,$\\
 & $\quad -4t^{3}-8t^{2}+4t-4\rangle$\\
\addlinespace[2pt]
\multicolumn{2}{@{}p{0.97\textwidth}@{}}{$i=23$, $D_i=52$, $\rho_{ij}$: $3{:}2, \allowbreak 4{:}2, \allowbreak 5{:}1, \allowbreak 12{:}2, \allowbreak 18{:}1, \allowbreak 26{:}2, \allowbreak 28{:}2$}\\
$w_i$ & $\langle -316t^{3}-148t^{2}+196t+124,$\\
 & $\quad -120t^{3}+188t^{2}+160t+80\rangle$\\
\addlinespace[2pt]
\multicolumn{2}{@{}p{0.97\textwidth}@{}}{$i=24$, $D_i=4$, $\rho_{ij}$: $5{:}2, \allowbreak 11{:}1, \allowbreak 18{:}2, \allowbreak 28{:}1, \allowbreak 31{:}2$}\\
$w_i$ & $\langle 4t^{3}-4t^{2},$\\
 & $\quad -4t^{5}+4t^{4}+8t^{3}-8t^{2}+4\rangle$\\
\addlinespace[2pt]
\multicolumn{2}{@{}p{0.97\textwidth}@{}}{$i=25$, $D_i=4$, $\rho_{ij}$: $3{:}1, \allowbreak 5{:}2, \allowbreak 7{:}2, \allowbreak 12{:}1, \allowbreak 13{:}2, \allowbreak 16{:}2, \allowbreak 30{:}2$}\\
$w_i$ & $\langle 10t^{3}+2t^{2}+4t+8,$\\
 & $\quad -18t^{3}-12t^{2}+6\rangle$\\
\addlinespace[2pt]
\multicolumn{2}{@{}p{0.97\textwidth}@{}}{$i=26$, $D_i=324$, $\rho_{ij}$: $0{:}2, \allowbreak 2{:}2, \allowbreak 12{:}1, \allowbreak 13{:}1, \allowbreak 14{:}2, \allowbreak 16{:}2, \allowbreak 18{:}2, \allowbreak 21{:}1, \allowbreak 23{:}1, \allowbreak 28{:}1, \allowbreak 30{:}1, \allowbreak 31{:}2$}\\
$w_i$ & $\langle -6804t^{5}-5832t^{4}-1620t^{3}+648t^{2}-1296t-648,$\\
 & $\quad -5508t^{5}-22032t^{4}-11340t^{3}-648t^{2}+2268t-5508\rangle$\\
\addlinespace[2pt]
\multicolumn{2}{@{}p{0.97\textwidth}@{}}{$i=27$, $D_i=1$, $\rho_{ij}$: $6{:}1, \allowbreak 10{:}2, \allowbreak 11{:}2, \allowbreak 15{:}2, \allowbreak 16{:}1, \allowbreak 28{:}1$}\\
$w_i$ & $\langle 4t^{3}+6t^{2}+4t+2,$\\
 & $\quad 4t^{3}+8t^{2}+6t+2\rangle$\\
\addlinespace[2pt]
\multicolumn{2}{@{}p{0.97\textwidth}@{}}{$i=28$, $D_i=2$, $\rho_{ij}$: $16{:}1, \allowbreak 22{:}1, \allowbreak 23{:}1, \allowbreak 24{:}2, \allowbreak 26{:}2, \allowbreak 27{:}2$}\\
$w_i$ & $\langle 16t^{3}+20t^{2}+4t,$\\
 & $\quad 8t^{3}+16t^{2}+8t\rangle$\\
\addlinespace[2pt]
\multicolumn{2}{@{}p{0.97\textwidth}@{}}{$i=29$, $D_i=2$, $\rho_{ij}$: $2{:}1, \allowbreak 4{:}1, \allowbreak 7{:}2, \allowbreak 8{:}2, \allowbreak 16{:}1$}\\
$w_i$ & $\langle 2,$\\
 & $\quad -2t+2\rangle$\\
\addlinespace[2pt]
\multicolumn{2}{@{}p{0.97\textwidth}@{}}{$i=30$, $D_i=10$, $\rho_{ij}$: $6{:}1, \allowbreak 10{:}1, \allowbreak 12{:}1, \allowbreak 17{:}1, \allowbreak 18{:}1, \allowbreak 25{:}1, \allowbreak 26{:}2$}\\
$w_i$ & $\langle 8t^{5}+24t^{4}-20t^{3}+32t^{2}+40t+24,$\\
 & $\quad 4t^{5}+32t^{4}+16t^{2}+20t+12\rangle$\\
\addlinespace[2pt]
\multicolumn{2}{@{}p{0.97\textwidth}@{}}{$i=31$, $D_i=2$, $\rho_{ij}$: $5{:}2, \allowbreak 21{:}2, \allowbreak 24{:}1, \allowbreak 26{:}1$}\\
$w_i$ & $\langle -2t^{3}+4t-2,$\\
 & $\quad 2t^{3}-6t^{2}+8t-4\rangle$\\
\addlinespace[2pt]
\end{longtable}
\endgroup

\begingroup\small
\setlength{\LTcapwidth}{\textwidth}
\begin{longtable}{@{}lrl@{\qquad}lrl@{}}
\caption{The edges $(i,j)$ with a linear endpoint, the linear endpoint $k$ used, and the norm of $R_{ij}=\operatorname{Res}_t(f_k,g_{kl})$.}\label{tab:sep}\\
\hline $(i,j)$ & $k$ & $N(R_{ij})$ & $(i,j)$ & $k$ & $N(R_{ij})$\\\hline\endfirsthead\hline $(i,j)$ & $k$ & $N(R_{ij})$ & $(i,j)$ & $k$ & $N(R_{ij})$\\\hline\endhead\hline\endfoot
\multicolumn{4}{@{}l}{\emph{The first certificate.}}\\
$(0,1)$ & 0 & $2^{4}$ & $(7,9)$ & 7 & $7$\\
$(0,4)$ & 0 & $2^{4}$ & $(7,14)$ & 7 & $2^{2}\cdot 7$\\
$(0,5)$ & 0 & $2^{4}$ & $(7,19)$ & 7 & $1$\\
$(0,7)$ & 0 & $2^{2}$ & $(8,9)$ & 8 & $1$\\
$(0,10)$ & 0 & $2^{4}$ & $(8,15)$ & 8 & $2^{4}$\\
$(0,13)$ & 0 & $1$ & $(8,16)$ & 8 & $3$\\
$(0,15)$ & 0 & $2^{2}$ & $(9,14)$ & 9 & $7^{3}$\\
$(0,16)$ & 0 & $2^{2}\cdot 3^{2}$ & $(9,20)$ & 9 & $3$\\
$(0,18)$ & 0 & $2^{8}\cdot 3^{2}$ & $(9,22)$ & 9 & $1$\\
$(0,27)$ & 0 & $1$ & $(10,20)$ & 10 & $7$\\
$(1,2)$ & 1 & $3$ & $(10,22)$ & 10 & $2^{2}$\\
$(1,3)$ & 1 & $1$ & $(10,25)$ & 10 & $2^{4}\cdot 7$\\
$(1,5)$ & 1 & $2^{2}\cdot 7$ & $(11,15)$ & 11 & $2^{4}$\\
$(1,8)$ & 1 & $2^{8}\cdot 3$ & $(11,16)$ & 11 & $2^{2}\cdot 3$\\
$(1,11)$ & 1 & $2^{4}$ & $(11,17)$ & 11 & $1$\\
$(1,12)$ & 1 & $7$ & $(11,22)$ & 11 & $2^{2}$\\
$(1,15)$ & 1 & $2^{2}$ & $(11,27)$ & 11 & $1$\\
$(1,23)$ & 1 & $2^{4}$ & $(12,13)$ & 12 & $3^{3}\cdot 7$\\
$(2,5)$ & 2 & $1$ & $(12,14)$ & 12 & $2^{6}\cdot 7$\\
$(2,9)$ & 2 & $1$ & $(12,28)$ & 12 & $2^{2}$\\
$(2,16)$ & 2 & $3$ & $(13,14)$ & 13 & $2^{4}\cdot 7$\\
$(2,17)$ & 2 & $1$ & $(13,15)$ & 13 & $2^{4}$\\
$(3,12)$ & 3 & $1$ & $(13,16)$ & 13 & $3^{3}$\\
$(3,16)$ & 3 & $1$ & $(13,21)$ & 13 & $3$\\
$(3,18)$ & 3 & $1$ & $(13,27)$ & 13 & $3^{2}$\\
$(3,19)$ & 3 & $1$ & $(14,15)$ & 14 & $2^{8}$\\
$(3,23)$ & 3 & $1$ & $(15,19)$ & 15 & $2^{4}$\\
$(4,18)$ & 4 & $3\cdot 7$ & $(17,21)$ & 17 & $1$\\
$(4,23)$ & 4 & $2^{2}\cdot 3$ & $(17,27)$ & 17 & $1$\\
$(4,24)$ & 4 & $2^{2}\cdot 7$ & $(18,24)$ & 18 & $2^{6}\cdot 3^{2}\cdot 7$\\
$(4,27)$ & 4 & $1$ & $(18,26)$ & 18 & $3$\\
$(5,9)$ & 5 & $3\cdot 7$ & $(19,24)$ & 19 & $2^{4}\cdot 7$\\
$(5,19)$ & 5 & $1$ & $(19,26)$ & 19 & $7$\\
$(5,23)$ & 5 & $3$ & $(20,21)$ & 20 & $3^{5}\cdot 7$\\
$(6,9)$ & 6 & $3^{2}$ & $(21,25)$ & 21 & $7$\\
$(6,10)$ & 6 & $2^{4}$ & $(22,28)$ & 22 & $2^{2}$\\
$(6,11)$ & 6 & $2^{4}\cdot 3^{2}$ & $(23,26)$ & 23 & $1$\\
$(6,18)$ & 6 & $2^{4}\cdot 3^{2}$ & $(24,26)$ & 24 & $3^{2}\cdot 7$\\
$(6,21)$ & 6 & $3^{2}$ & $(24,27)$ & 24 & $3^{2}\cdot 7^{3}$\\
$(6,28)$ & 6 & $2^{4}$ & $(26,28)$ & 26 & $1$\\
\multicolumn{4}{@{}l}{\emph{The second certificate.}}\\
$(0,15)$ & 0 & $1$ & $(8,17)$ & 8 & $7$\\
$(0,26)$ & 0 & $2^{2}$ & $(8,19)$ & 8 & $2^{4}\cdot 7^{2}$\\
$(1,6)$ & 1 & $3^{2}$ & $(8,21)$ & 8 & $2^{2}\cdot 7\cdot 13$\\
$(1,11)$ & 1 & $3$ & $(8,22)$ & 8 & $5^{2}\cdot 7^{2}$\\
$(1,12)$ & 1 & $2^{2}\cdot 3$ & $(8,29)$ & 8 & $1$\\
$(1,22)$ & 1 & $3^{2}$ & $(9,21)$ & 9 & $1$\\
$(2,5)$ & 2 & $2^{2}\cdot 3^{2}$ & $(10,15)$ & 10 & $3$\\
$(2,6)$ & 2 & $3^{2}\cdot 13$ & $(10,27)$ & 10 & $2^{2}$\\
$(2,7)$ & 2 & $7\cdot 13$ & $(10,30)$ & 10 & $2^{2}\cdot 3$\\
$(2,8)$ & 2 & $7^{2}$ & $(11,18)$ & 11 & $2^{2}$\\
$(2,10)$ & 2 & $2^{2}\cdot 3$ & $(11,20)$ & 11 & $3^{2}$\\
$(2,11)$ & 2 & $3$ & $(11,21)$ & 11 & $2^{2}$\\
$(2,13)$ & 2 & $3^{2}$ & $(11,24)$ & 11 & $3$\\
$(2,17)$ & 2 & $3\cdot 7$ & $(11,27)$ & 11 & $1$\\
$(2,18)$ & 2 & $1$ & $(12,23)$ & 12 & $2^{4}$\\
$(2,19)$ & 2 & $3^{2}\cdot 7\cdot 13$ & $(12,25)$ & 12 & $1$\\
$(2,21)$ & 2 & $2^{2}\cdot 7\cdot 13$ & $(12,26)$ & 12 & $3$\\
$(2,26)$ & 2 & $3\cdot 7$ & $(12,30)$ & 12 & $2^{2}\cdot 3$\\
$(2,29)$ & 2 & $3$ & $(13,18)$ & 13 & $2^{4}$\\
$(3,4)$ & 3 & $2^{2}\cdot 3^{6}$ & $(13,21)$ & 13 & $2^{4}$\\
$(3,8)$ & 3 & $2^{4}\cdot 7\cdot 13$ & $(13,22)$ & 13 & $3^{2}$\\
$(3,16)$ & 3 & $2^{2}\cdot 3\cdot 13$ & $(13,25)$ & 13 & $2^{4}$\\
$(3,21)$ & 3 & $2^{2}\cdot 7\cdot 13^{2}$ & $(13,26)$ & 13 & $2^{4}\cdot 3^{2}$\\
$(3,23)$ & 3 & $3$ & $(14,15)$ & 14 & $1$\\
$(3,25)$ & 3 & $2^{8}$ & $(14,21)$ & 14 & $2^{2}$\\
$(4,5)$ & 4 & $2^{8}\cdot 3^{2}$ & $(14,22)$ & 14 & $1$\\
$(4,17)$ & 4 & $3\cdot 7$ & $(14,26)$ & 14 & $7^{2}$\\
$(4,20)$ & 4 & $2^{6}\cdot 3^{2}\cdot 7$ & $(15,22)$ & 15 & $2^{2}\cdot 3\cdot 5^{2}$\\
$(4,21)$ & 4 & $2^{2}\cdot 7\cdot 13^{2}$ & $(15,27)$ & 15 & $2^{2}$\\
$(4,22)$ & 4 & $2^{4}\cdot 3^{3}$ & $(16,21)$ & 16 & $2^{4}\cdot 13$\\
$(4,23)$ & 4 & $2^{2}\cdot 3\cdot 13$ & $(16,25)$ & 16 & $2^{4}\cdot 7$\\
$(4,29)$ & 4 & $3$ & $(16,26)$ & 16 & $7\cdot 13$\\
$(5,10)$ & 5 & $2^{4}$ & $(16,27)$ & 16 & $2^{4}$\\
$(5,14)$ & 5 & $2^{2}$ & $(16,28)$ & 16 & $3$\\
$(5,16)$ & 5 & $2^{4}$ & $(16,29)$ & 16 & $3$\\
$(5,18)$ & 5 & $2^{4}$ & $(17,19)$ & 17 & $3\cdot 7$\\
$(5,19)$ & 5 & $2^{4}$ & $(17,20)$ & 17 & $7^{3}$\\
$(5,23)$ & 5 & $2^{4}$ & $(17,21)$ & 17 & $2^{4}\cdot 7^{2}$\\
$(5,24)$ & 5 & $1$ & $(17,22)$ & 17 & $3\cdot 7$\\
$(5,25)$ & 5 & $2^{4}$ & $(17,30)$ & 17 & $3^{2}$\\
$(5,31)$ & 5 & $2^{4}$ & $(18,19)$ & 18 & $2^{10}$\\
$(6,12)$ & 6 & $1$ & $(18,23)$ & 18 & $2^{4}\cdot 7$\\
$(6,17)$ & 6 & $7$ & $(18,24)$ & 18 & $1$\\
$(6,19)$ & 6 & $13$ & $(18,26)$ & 18 & $2^{2}\cdot 7^{2}$\\
$(6,20)$ & 6 & $2^{4}\cdot 7$ & $(18,30)$ & 18 & $2^{4}\cdot 7$\\
$(6,27)$ & 6 & $2^{4}$ & $(19,21)$ & 19 & $2^{2}\cdot 7\cdot 13$\\
$(6,30)$ & 6 & $1$ & $(20,21)$ & 20 & $2^{2}\cdot 7$\\
$(7,8)$ & 7 & $5^{2}\cdot 13$ & $(21,26)$ & 21 & $2^{2}\cdot 13^{2}$\\
$(7,9)$ & 7 & $1$ & $(21,31)$ & 21 & $2^{4}$\\
$(7,14)$ & 7 & $7^{2}$ & $(22,28)$ & 22 & $2^{8}$\\
$(7,16)$ & 7 & $7\cdot 13$ & $(23,26)$ & 23 & $7\cdot 13$\\
$(7,17)$ & 7 & $1$ & $(23,28)$ & 23 & $2^{2}\cdot 3\cdot 13$\\
$(7,19)$ & 7 & $13^{2}$ & $(24,28)$ & 24 & $3^{4}$\\
$(7,22)$ & 7 & $5^{2}$ & $(24,31)$ & 24 & $3^{2}$\\
$(7,25)$ & 7 & $7$ & $(25,30)$ & 25 & $2^{4}\cdot 7$\\
$(7,29)$ & 7 & $1$ & $(26,28)$ & 26 & $3^{2}\cdot 13$\\
$(8,9)$ & 8 & $1$ & $(26,30)$ & 26 & $3^{4}\cdot 7^{2}$\\
$(8,13)$ & 8 & $2^{4}$ & $(26,31)$ & 26 & $2^{2}\cdot 3^{2}$\\
$(8,15)$ & 8 & $5^{2}$ & $(27,28)$ & 27 & $2^{6}$\\
\end{longtable}
\endgroup

\begingroup\small
\setlength{\LTcapwidth}{\textwidth}
\begin{longtable}{@{}rrl>{\raggedright\arraybackslash}p{0.66\textwidth}@{}}
\caption{The terminal residue classes $x\equiv a$, $y\equiv b\pmod{p^k}$ of Step~4, grouped by $p$, $k$ and the value of $(B^{(g)}_p)_g$.}\label{tab:local}\\
\hline $p$ & $k$ & value & classes $(a,b)$\\\hline\endfirsthead
\hline $p$ & $k$ & value & classes $(a,b)$\\\hline\endhead\hline\endfoot
2 & 3 & $(1,2)$ & $(5,0)$,\allowbreak\ $(5,2)$,\allowbreak\ $(5,4)$,\allowbreak\ $(5,6)$\\
\addlinespace[3pt]
3 & 3 & $(0,1)$ & $(22,2)$,\allowbreak\ $(22,11)$,\allowbreak\ $(22,20)$\\
3 & 3 & $(0,2)$ & $(5,17)$,\allowbreak\ $(14,8)$,\allowbreak\ $(23,26)$\\
3 & 3 & $(1,1)$ & $(2,11)$,\allowbreak\ $(11,2)$,\allowbreak\ $(20,20)$\\
3 & 3 & $(1,2)$ & $(4,8)$,\allowbreak\ $(4,17)$,\allowbreak\ $(4,26)$\\
3 & 3 & $(2,0)$ & $(8,14)$,\allowbreak\ $(13,5)$,\allowbreak\ $(13,14)$,\allowbreak\ $(13,23)$,\allowbreak\ $(17,5)$,\allowbreak\ $(26,23)$\\
\addlinespace[3pt]
5 & 1 & $(0,0)$ & $(0,3)$,\allowbreak\ $(1,1)$,\allowbreak\ $(2,0)$,\allowbreak\ $(2,2)$,\allowbreak\ $(2,3)$\\
\addlinespace[3pt]
7 & 1 & $(0,0)$ & $(0,2)$\\
7 & 2 & $(0,0)$ & $(2,14)$,\allowbreak\ $(3,11)$,\allowbreak\ $(5,45)$,\allowbreak\ $(9,0)$,\allowbreak\ $(10,32)$,\allowbreak\ $(16,35)$,\allowbreak\ $(17,4)$,\allowbreak\ $(18,10)$,\allowbreak\ $(19,31)$,\allowbreak\ $(23,21)$,\allowbreak\ $(24,25)$,\allowbreak\ $(25,17)$,\allowbreak\ $(26,24)$,\allowbreak\ $(30,7)$,\allowbreak\ $(31,19)$,\allowbreak\ $(31,33)$,\allowbreak\ $(31,46)$,\allowbreak\ $(32,24)$,\allowbreak\ $(33,17)$,\allowbreak\ $(37,42)$,\allowbreak\ $(38,18)$,\allowbreak\ $(39,31)$,\allowbreak\ $(40,10)$,\allowbreak\ $(44,28)$,\allowbreak\ $(45,39)$,\allowbreak\ $(46,38)$,\allowbreak\ $(47,3)$\\
7 & 3 & $(0,0)$ & $(4,290)$,\allowbreak\ $(11,248)$,\allowbreak\ $(12,136)$,\allowbreak\ $(53,339)$,\allowbreak\ $(60,297)$,\allowbreak\ $(61,87)$,\allowbreak\ $(102,45)$,\allowbreak\ $(109,3)$,\allowbreak\ $(110,38)$,\allowbreak\ $(129,26)$,\allowbreak\ $(129,75)$,\allowbreak\ $(129,124)$,\allowbreak\ $(129,173)$,\allowbreak\ $(129,222)$,\allowbreak\ $(129,271)$,\allowbreak\ $(129,320)$,\allowbreak\ $(151,94)$,\allowbreak\ $(158,52)$,\allowbreak\ $(159,332)$,\allowbreak\ $(178,5)$,\allowbreak\ $(178,47)$,\allowbreak\ $(178,54)$,\allowbreak\ $(178,96)$,\allowbreak\ $(178,103)$,\allowbreak\ $(178,145)$,\allowbreak\ $(178,152)$,\allowbreak\ $(178,194)$,\allowbreak\ $(178,201)$,\allowbreak\ $(178,243)$,\allowbreak\ $(178,250)$,\allowbreak\ $(178,292)$,\allowbreak\ $(178,299)$,\allowbreak\ $(178,341)$,\allowbreak\ $(200,143)$,\allowbreak\ $(207,101)$,\allowbreak\ $(208,283)$,\allowbreak\ $(227,12)$,\allowbreak\ $(227,40)$,\allowbreak\ $(227,61)$,\allowbreak\ $(227,89)$,\allowbreak\ $(227,110)$,\allowbreak\ $(227,138)$,\allowbreak\ $(227,159)$,\allowbreak\ $(227,187)$,\allowbreak\ $(227,208)$,\allowbreak\ $(227,236)$,\allowbreak\ $(227,257)$,\allowbreak\ $(227,285)$,\allowbreak\ $(227,306)$,\allowbreak\ $(227,334)$,\allowbreak\ $(249,192)$,\allowbreak\ $(256,150)$,\allowbreak\ $(257,234)$,\allowbreak\ $(298,241)$,\allowbreak\ $(305,199)$,\allowbreak\ $(306,185)$\\
\addlinespace[3pt]
13 & 1 & $(0,0)$ & $(2,3)$,\allowbreak\ $(3,1)$,\allowbreak\ $(4,8)$,\allowbreak\ $(4,10)$,\allowbreak\ $(5,1)$,\allowbreak\ $(6,7)$,\allowbreak\ $(7,0)$,\allowbreak\ $(9,2)$,\allowbreak\ $(10,4)$,\allowbreak\ $(12,5)$\\
13 & 2 & $(0,0)$ & $(8,163)$,\allowbreak\ $(11,166)$,\allowbreak\ $(21,137)$,\allowbreak\ $(24,75)$,\allowbreak\ $(34,111)$,\allowbreak\ $(37,153)$,\allowbreak\ $(47,85)$,\allowbreak\ $(50,62)$,\allowbreak\ $(63,140)$,\allowbreak\ $(73,33)$,\allowbreak\ $(76,49)$,\allowbreak\ $(86,7)$,\allowbreak\ $(99,150)$,\allowbreak\ $(102,36)$,\allowbreak\ $(112,124)$,\allowbreak\ $(115,114)$,\allowbreak\ $(125,98)$,\allowbreak\ $(128,23)$,\allowbreak\ $(138,72)$,\allowbreak\ $(141,101)$,\allowbreak\ $(151,46)$,\allowbreak\ $(154,10)$,\allowbreak\ $(164,20)$,\allowbreak\ $(167,88)$\\
13 & 3 & $(0,0)$ & $(60,59)$,\allowbreak\ $(89,972)$,\allowbreak\ $(229,1918)$,\allowbreak\ $(258,1986)$,\allowbreak\ $(398,1580)$,\allowbreak\ $(427,803)$,\allowbreak\ $(567,1242)$,\allowbreak\ $(596,1817)$,\allowbreak\ $(736,904)$,\allowbreak\ $(765,634)$,\allowbreak\ $(905,566)$,\allowbreak\ $(934,1648)$,\allowbreak\ $(1074,228)$,\allowbreak\ $(1103,465)$,\allowbreak\ $(1243,2087)$,\allowbreak\ $(1272,1479)$,\allowbreak\ $(1412,1749)$,\allowbreak\ $(1441,296)$,\allowbreak\ $(1581,1411)$,\allowbreak\ $(1610,1310)$,\allowbreak\ $(1750,1073)$,\allowbreak\ $(1779,127)$,\allowbreak\ $(1919,735)$,\allowbreak\ $(1948,1141)$,\allowbreak\ $(2088,397)$,\allowbreak\ $(2117,2155)$\\
\addlinespace[3pt]
\end{longtable}
\endgroup
\end{document}
